\documentclass[10pt,a4paper]{article}

\usepackage[T1]{fontenc}
\usepackage[utf8]{inputenc}
\usepackage[french]{babel}
\usepackage{lmodern}
\usepackage{microtype}
\usepackage[a4paper,margin=2.2cm]{geometry}
\usepackage{enumitem}
\usepackage[hidelinks]{hyperref}

\setlength{\parindent}{0pt}
\setlength{\parskip}{0.6em}
\setlist[itemize]{leftmargin=1.5em,itemsep=0.4em,topsep=0.4em}

\newcommand{\projectversion}{1}
\newcommand{\projectdate}{Version initiale transcrite le 12 juillet 2026}

\title{L'émergence de l'intéressant\\
\large Naissance des idées, compréhension et singularité des formes}
\author{François Pachet}
\date{Projet de thèse -- version \projectversion\\\small \projectdate}

\begin{document}

\maketitle

\begin{center}
\emph{Chercheur en IA (thèse, HDR, ancien directeur de recherche dans le privé,
Sony puis Spotify)}
\end{center}

Ce thème et ces intuitions sont fondés sur mon expérience de chercheur en IA,
caractérisée par une approche constructiviste et empirique. J'ai passé quarante
ans à construire des systèmes d'IA pour « générer » des musiques et des textes,
avec obsession et parfois avec succès. J'aimerais maintenant essayer de
comprendre ce que je fais.

\section*{Projet de thèse}

Ce projet de thèse a pour objet une question à la fois simple et évanescente :
qu'est-ce qui fait qu'une idée, une œuvre, une phrase, une mélodie ou un objet
nous paraît intéressant ? Cette question engage plusieurs problèmes que la
philosophie a abordés dans des cadres distincts : la naissance des idées, depuis
les analyses de l'accès à l'intelligible chez Platon et Aristote jusqu'aux
réflexions sur l'invention, l'abduction en logique ou la découverte chez Peirce,
Poincaré ou Hadamard ; l'expérience de la compréhension, de la théorie
aristotélicienne du savoir aux analyses du jugement et du sens chez Kant, Husserl
ou Wittgenstein ; la soudaineté de certaines évidences, depuis l'intuition
cartésienne ou la \emph{scientia intuitiva} de Spinoza jusqu'aux changements
d'aspect analysés par Wittgenstein ; enfin le statut de la mémoire, de
l'attention et de la valeur, particulièrement chez Bergson, Merleau-Ponty,
Goodman ou Deleuze.

On peut partir d'une série d'expériences très communes. Il arrive qu'une idée
surgisse « d'un coup », qu'une pensée jusque-là diffuse se cristallise soudain,
qu'un problème devienne compréhensible, qu'une formule s'impose avec une forme
d'évidence. Il arrive aussi qu'un objet nous retienne de manière
disproportionnée : une chanson, une modulation, un passage littéraire, une image,
un raisonnement nous fascinent, au point d'appeler la reprise, la réécoute,
l'interprétation, l'analyse. Or cette fascination ne se laisse réduire ni à la
nouveauté brute, ni à la beauté au sens classique, ni à la préférence subjective,
ni à une propriété objective isolable dans l'objet. Elle semble tenir à une
relation plus complexe entre une forme, une mémoire, une artillerie d'attentes et
une capacité de reconnaissance.

Le projet défend l'hypothèse que ces phénomènes relèvent d'une même question :
celle de l'émergence de l'intéressant. L'intéressant n'est pas compris ici comme
un qualificatif mondain, mais comme une possible catégorie philosophique,
désignant le moment où quelque chose se détache pour un esprit comme une forme,
au sens fort d'une organisation saillante et structurée, sens que ce terme peut
recevoir aussi bien dans la tradition aristotélicienne ou phénoménologique que,
plus récemment, chez Jean Petitot. L'ambition de cette thèse sera ainsi de
constituer une philosophie de l'intéressant, en montrant que la naissance des
idées, la compréhension et la fascination pour certaines formes ne sont pas trois
problèmes indépendants, mais trois aspects d'un même phénomène.

La question directrice peut être formulée ainsi : sous quelles conditions quelque
chose apparaît-il à un sujet comme intéressant ? Cette question en entraîne
plusieurs autres. L'intéressant est-il dans l'objet lui-même, dans le sujet qui le
reçoit, ou dans une relation entre les deux, difficulté déjà au cœur de l'analyse
kantienne du jugement de goût, qui ne relève ni d'une propriété objective des
choses ni d'une simple préférence privée ? Une chanson fascinante possède-t-elle
des propriétés propres qui expliquent sa puissance, selon une hypothèse proche
des ambitions du \emph{Hit Song Science}, ou bien cette puissance dépend-elle
essentiellement de la mémoire, de l'histoire, de la culture et des attentes du
sujet ? Peut-on penser l'intéressant sans le réduire ni à l'objectivité brute, ni
au subjectivisme du goût ?

Le projet fait l'hypothèse que l'intéressant désigne une relation de saillance
(d'attention ?) entre une forme et un esprit historiquement, affectivement et
culturellement constitué. Une chose apparaît intéressante lorsqu'elle réalise
plusieurs conditions à la fois : une certaine cohérence interne, sans quoi elle
reste confuse ; une tension, une surprise, sans quoi elle se dissout dans le
déjà-vu ; enfin une puissance de relance, c'est-à-dire la capacité d'ouvrir un
espace de reprises, de variations et de pensées possibles. L'intéressant serait
ainsi ce qui conjugue intelligibilité, singularité et fécondité.

Cette hypothèse sera poussée plus loin. Certains objets particulièrement
intéressants se présentent comme des singularités de forme : ils semblent
résoudre de manière nécessaire un problème qui n'était pas pleinement formulé
avant eux. Ils ne se contentent pas de répondre à une question donnée ; ils font
apparaître, par leur existence même, la question à laquelle ils répondent. En ce
sens, l'objet intéressant comme nécessité inventée peut être rapproché de Valéry,
lorsqu'il décrit l'invention comme production d'une forme qui découvre ses
propres contraintes, ou de Bachelard, lorsqu'il montre que certaines avancées de
la pensée reconfigurent le champ même des problèmes.

L'un des enjeux de la recherche sera de clarifier le rapport entre objet, mémoire
et sujet dans l'expérience de l'intéressant. Une première tentation consisterait
à localiser l'intéressant dans l'objet ; une seconde à tout rabattre sur le sujet.
Le projet cherchera à dépasser cette opposition. Il fera l'hypothèse que
l'intéressant est une propriété relationnelle au sens fort : il dépend d'une
histoire du sujet, de ses attentes, de sa mémoire, de sa culture et de ses
aptitudes perceptives ; mais il suppose aussi dans l'objet une organisation, une
économie, une structure de tensions et de résolutions. Ce qui intéresse, c'est
une rencontre entre une forme et une mémoire capable d'en éprouver la
singularité. Sur ce point, le projet pourra convoquer les analyses bergsoniennes
de la mémoire et de la durée, les descriptions husserliennes de la rétention et
de la protention, ou encore la pensée merleau-pontienne du sens comme émergence
dans un champ.

La musique constitue ici un terrain privilégié. En tant qu'objet temporel, elle
met à nu les rapports entre mémoire, attente, surprise et reconnaissance.
Certaines chansons fascinent parce qu'elles semblent installer leur propre
espace de contraintes, puis résoudre les attentes qu'elles ont fait naître d'une
manière qui paraît après coup évidente mais qu'il était impossible d'anticiper.
Une attention particulière sera portée à certaines chansons relevant de ce que
je propose d'appeler, de manière encore exploratoire, le régime Jotney,
c'est-à-dire des formes où la surprise harmonique et mélodique s'accompagne d'une
impression remarquable de nécessité locale. La musique fonctionnera ainsi non
comme simple cas particulier, mais comme laboratoire philosophique de
l'intéressant.

Un autre axe important sera la relation entre intéressant, rareté et difficulté
de production. Il ne s'agit pas seulement de dire que l'intéressant est le rare
au sens statistique, mais qu'il existe des formes dont on sent qu'on ne pourrait
déplacer un élément sans altérer l'ensemble (autostables). À cela s'ajoute parfois
une impression de simplicité de surface qui dissimule une grande difficulté
d'invention. Ce phénomène, que l'on peut relier à la notion de difficultuosité
(développée dans un autre contexte), pourra être reformulé philosophiquement :
certaines formes nous frappent parce qu'elles paraissent aller de soi, tout en
laissant pressentir l'improbabilité de leur genèse.

La méthode : analyse conceptuelle, étude d'auteurs, confrontation de modèles de
l'intelligibilité, de l'invention et de la valeur. Elle s'appuiera sur des cas
pris dans la philosophie, la littérature, les mathématiques et la musique. Le
projet pourra aussi convoquer certains travaux contemporains en sciences
cognitives, en psychologie de l'\emph{insight}, en théorie de la perception ou en
modélisation de l'attention, sans pour autant naturaliser la question. Un ancien
projet consacré à l'\emph{interestingness}, élaboré dans un cadre computationnel,
pourra jouer ici un rôle d'arrière-plan conceptuel, à condition que ses intuitions
soient reprises et reformulées dans un cadre proprement philosophique.

L'organisation provisoire de la recherche pourrait suivre cinq phases : une
phénoménologie de l'intéressant ; une analyse de la compréhension comme émergence
de l'intelligible ; une réflexion sur la naissance des idées ; une théorie de la
singularité intéressante ; enfin une mise à l'épreuve de ces hypothèses dans le
cas de la musique.

L'originalité de cette thèse tient à la constitution d'un objet philosophique
rarement traité comme tel : l'intéressant. Là où la philosophie a souvent
privilégié le vrai, le beau, le bien, le sens ou le jugement, il s'agira ici de
prendre au sérieux cette catégorie intermédiaire et décisive, qui semble gouverner
une large part de notre vie intellectuelle et esthétique, voire du quotidien.
L'enjeu est de proposer une conception de la pensée qui ne la réduise ni à
l'application de règles, ni à la simple réception de contenus, ni à la
reconstruction logique après coup. La philosophie de l'intéressant serait alors
une manière renouvelée de penser ensemble l'invention, la compréhension et la
valeur.

\section*{Bibliographie initiale}

\begin{itemize}
  \item Platon, \emph{La République} ; \emph{Théétète}.
  \item Aristote, \emph{Métaphysique} ; \emph{Seconds Analytiques} ;
        \emph{Poétique}.
  \item Descartes, \emph{Règles pour la direction de l'esprit} ;
        \emph{Méditations métaphysiques}.
  \item Spinoza, \emph{Éthique}.
  \item Kant, \emph{Critique de la raison pure} ;
        \emph{Critique de la faculté de juger}.
  \item Bergson, \emph{L'évolution créatrice} ; \emph{La pensée et le mouvant}.
  \item Husserl, \emph{Expérience et jugement} ;
        \emph{Logique formelle et logique transcendantale}.
  \item Merleau-Ponty, \emph{Phénoménologie de la perception}.
  \item Wittgenstein, \emph{Recherches philosophiques}.
  \item Peirce, textes sur l'abduction et la formation des hypothèses.
  \item Bachelard, \emph{La formation de l'esprit scientifique}.
  \item Valéry, \emph{Cahiers} ; \emph{Tel Quel}.
  \item Goodman, \emph{Languages of Art}.
  \item Deleuze, \emph{Différence et répétition}.
  \item Hadamard, \emph{Essai sur la psychologie de l'invention dans le domaine
        mathématique}.
  \item Poincaré, \emph{Science et méthode}.
  \item Koestler, \emph{The Act of Creation}.
  \item Berlyne, \emph{Studies in the New Experimental Aesthetics}.
  \item Schmidhuber, « What's Interesting? »
  \item Abdallah et Plumbley, travaux sur attente et surprise en musique.
\end{itemize}

\section*{Travaux personnels en lien avec le sujet}

Liste complète de publications :
\url{https://www.francoispachet.fr/publications}

\begin{itemize}
  \item Pachet, F. (2026), \emph{La virtuosité à la portée des caniches}
        (en cours de publication). Ce texte porte sur les différentes formes de
        virtuosité et leurs effets recherchés sur les auditeurs.
  \item Pachet, F. (2026), \emph{De l'impossibilité de créer quoi que ce soit}
        (en cours de publication). Cet essai recense toutes les raisons
        scientifiques s'opposant à l'idée même de créer du nouveau.
  \item Pachet, F., Roy, P. et Carré, B. (2021), « Assisted Music Creation with
        Flow Machines: Towards New Categories of New », chapitre 18 de
        \emph{Handbook of Artificial Intelligence for Music}, Eduardo Miranda
        (dir.), Springer. J'y décris les résultats principaux du projet Flow
        Machines et conclus avec l'idée que l'IA renouvelle les catégories du
        nouveau.
  \item Pachet, F. (2018), \emph{Histoire d'une oreille}, Buchet-Chastel. Ce
        livre, augmenté avec des QR codes Spotify, est une tentative de
        description de l'ontogenèse d'une oreille musicale. Titre alternatif :
        \emph{Proust écoute les Beatles}.
  \item Pachet, F. (2012), « Musical Virtuosity and Creativity », dans
        J. McCormack et M. d'Inverno (dir.), \emph{Computers and Creativity},
        Springer. Je présente la virtuosité comme un trait important des
        productions artistiques et décris un système (Virtuoso) permettant de
        générer des phrases virtuoses en bebop en temps réel.
  \item Pachet, F. (2008), ancien projet ERC sur l'\emph{interestingness} et
        travaux connexes sur la singularité, la difficulté de production et la
        relation entre création et perception.
  \item Pachet, F. (2008), « Des machines à sortir de soi », dans Bernard
        Stiegler (dir.), \emph{Le design de nos existences à l'époque de
        l'innovation ascendante}, p. 229--251, Mille et une nuits, Fayard. J'y
        décris les Flow Machines et l'idée de construire des systèmes favorisant
        les états de flow selon Csíkszentmihályi.
  \item Pachet, F. (2006), « Les clefs d'une mélodie intéressante »,
        \emph{La Naissance de l'Art}, p. 193--201, Tallandier. Ce papier propose
        déjà diverses approches pour modéliser l'intéressant dans les mélodies.
\end{itemize}

\end{document}
